\paragraph{IV 20.13: Morphismes élémentaires et théorème de M. Artin (cote 25, sites 75--83).}
Definition 13.1 introduces elementary morphisms by removing an étale closed subscheme from a smooth projective family of geometrically connected curves, and calls composites of elementary morphisms poly-elementary. Theorem 13.2, with Artin's name and the hypotheses ``géom. irréd.'' and ``parfait infini'' added above the line at site 75, gives poly-elementary open neighbourhoods of a smooth point. Its induction begins with the normal projective model of a function field and explicitly refers to Chapter II, paragraph 7; the text also explains its use of the composite of an empty family. The reduction to a normal projective scheme, the singular loci of $X$ and $Y$, and the linear-section lemma 13.3 occupy sites 76--77. The lemma invokes the Veronese multiples from section 9, while its application uses the Grassmannian incidence notation and Bertini's connectedness statement to construct the elementary fibration (sites 77--78). The proof of the lemma leads to substantial ``remords'' to sections 8 and 9: site 79 cites 2.3 as revised in section 12 and rewrites the linear spaces through a fixed subspace as spaces through a point; sites 79--80 use 8.8 and the symmetry expressed in 8.7(c), then explain the induction on the codimension of a linear section. The handwritten cancellation at site 77 includes the entire clause asserting $C\cap Y=\emptyset$, while the added finite-set avoidance condition remains; at site 78 a whole typed interlinear line belongs after the incidence scheme $\underline{T}^{n-1}_Q$. Remark 13.4 discusses the field hypotheses, records a finite-extension conclusion, treats the omission of geometric irreducibility, and proposes the general-base variant through 7.1 (sites 81--82). Its inserted label c), its renumbering of the following paragraph to d), and the later typed c) and e) show the changing arrangement of these remarks. The finite-set refinement at site 82 and the discussion of the analytic space at site 83 complete the section; the latter describes the homotopy groups, a classifying space $K(\pi_1,1)$, and its universal covering space. Several printed expressions remain especially useful to check beside the scans: definition 13.1 has $X'\to X$ at site 75; the neighbourhood $V$ is said to lie ``dans X'' at site 78; site 82 first gives $L_i$ codimension $n-1$ and later writes codimension $n$. The manuscript's own qualifications and unresolved letterforms remain marked, and these expressions have received no mathematical repair.
